%% =====================================================================
%% SAGE-MM --- ACM Transactions on Embedded Computing Systems (TECS)
%% Journal submission manuscript.
%%
%% Template: ACM `acmart' document class, `acmsmall' journal format
%% (the format used by ACM TECS). See paper/README.md for build and
%% current publication-cost verification notes.
%%
%% Cost note: ACM announced a fully open-access publishing model beginning
%% in 2026. Author cost depends on ACM Open institutional coverage, waivers,
%% and the current ACM policy; this source cannot guarantee a zero APC.
%% =====================================================================
\documentclass[acmsmall,screen]{acmart}

%% ---------------------------------------------------------------------
%% Journal identification (ACM fills the final values at acceptance).
%% ---------------------------------------------------------------------
\acmJournal{TECS}
\acmVolume{0}
\acmNumber{0}
\acmArticle{0}
\acmMonth{1}
\acmYear{2026}
\acmDOI{}

%% Rights metadata is provisional; ACM supplies final values at acceptance.
\setcopyright{rightsretained}
\copyrightyear{2026}

%% Remove the "submitted to ACM" watermark line for a clean draft PDF.
\settopmatter{printacmref=true}

%% ---------------------------------------------------------------------
%% Packages
%% ---------------------------------------------------------------------
\usepackage{booktabs}
\usepackage{array}
%% Pseudocode for the controller decision procedure (Algorithm 1). These are
%% part of a standard ACM/texlive install; a minimal TeX environment that lacks
%% algorithmicx can comment them out together with the algorithm float below.
\usepackage{algorithm}
\usepackage{algpseudocode}
%% Keep builds deterministic on minimal TeX installations that lack ACM's
%% preferred scalable Libertine/NewTX fonts. Full ACM environments may remove
%% this line to restore microtype font expansion.
\microtypesetup{expansion=false}
\setlength{\headheight}{19pt}
%% Note: acmart already loads amsmath and the newtxmath symbol set; do not
%% re-load amssymb (it clashes with newtxmath on \Bbbk).

%% Internal layout mode. A submission build MUST set this false and provide
%% provenance-backed observed results. scripts/build_paper.sh --submission
%% refuses to build while this flag is true.
\newif\ifincludesimulation
\includesimulationtrue

\begin{document}

%% =====================================================================
%% Title and authors
%% =====================================================================
\title{SAGE-MM: Coordinating Heap Configuration, Interop Allocation, and
Page Reclamation in Memory-Constrained Embedded .NET Firmware}

\author{Geunsik Lim}
\authornote{Corresponding author.}
\email{leemgs@gmail.com}
\affiliation{%
  \institution{[Institution]}
  \city{[City]}
  \country{[Country]}
}

\renewcommand{\shortauthors}{Lim et al.}

%% =====================================================================
%% Abstract
%% =====================================================================
\begin{abstract}
Memory-constrained embedded devices co-locate a managed heap, native
interoperability, and file-backed executable mappings inside a single
process budget, yet these layers are conventionally tuned in isolation.
We present SAGE-MM, an engineering design that \emph{coordinates} three
deliberately distinct interventions under one embedded-process memory
budget: a static, architecture-specific managed-heap build configuration;
a source-reviewed value-type interop conversion that removes allocation;
and a bounded, online-controlled reclamation of clean file-backed pages
together with runtime compaction gating. SAGE-MM does not introduce a new
garbage collector, a new interop representation, a new reclamation
syscall, or a new learning algorithm. Its contribution is an implementable
cross-layer policy with an explicit action boundary, fail-closed safety
guards, and a falsifiable evaluation protocol for embedded .NET firmware.
The online controller consumes normalized garbage-collection pause,
fragmentation, page-fault, and resident-growth telemetry and selects a
reclamation interval using one of three interchangeable policies---a tuned
threshold rule, an exponentially weighted moving average (EWMA), or online
ridge regression---while a two-threshold fragmentation guard and a
deferral cap prevent compaction starvation. Invalid telemetry disables
reclamation and restores compaction without updating the model. We specify
a full-factorial evaluation over heap build, interop, reclamation, and
control, on the evaluated ARM32/ARM64 digital-television firmware plus at
least one independently assembled constrained Linux platform, reporting
per-mechanism and interaction effects on peak proportional set size (PSS),
tail pause, page faults, and interaction latency with run-level confidence
intervals. All reported effect sizes are generated from a single versioned
observed-result bundle; unmeasured quantities in this draft are marked and
must not be quoted as achieved improvements. The evidence is scoped to the
tested firmware and workloads; generality to other embedded managed
runtimes is stated as future validation.
\end{abstract}

%% ---------------------------------------------------------------------
%% ACM CCS concepts and keywords
%% ---------------------------------------------------------------------
\begin{CCSXML}
<ccs2012>
<concept>
<concept_id>10010520.10010521.10010542.10010544</concept_id>
<concept_desc>Computer systems organization~Embedded software</concept_desc>
<concept_significance>500</concept_significance>
</concept>
<concept>
<concept_id>10011007.10011074.10011099.10011692</concept_id>
<concept_desc>Software and its engineering~Garbage collection</concept_desc>
<concept_significance>500</concept_significance>
</concept>
<concept>
<concept_id>10011007.10010940.10010992.10010998</concept_id>
<concept_desc>Software and its engineering~Memory management</concept_desc>
<concept_significance>500</concept_significance>
</concept>
<concept>
<concept_id>10010583.10010588.10010559</concept_id>
<concept_desc>Hardware~Memory and dense storage</concept_desc>
<concept_significance>100</concept_significance>
</concept>
</ccs2012>
\end{CCSXML}

\ccsdesc[500]{Computer systems organization~Embedded software}
\ccsdesc[500]{Software and its engineering~Memory management}
\ccsdesc[500]{Software and its engineering~Garbage collection}

\keywords{Embedded runtimes, memory management, garbage collection,
managed/native interop, page reclamation, \texttt{madvise}, online control,
cross-layer coordination, .NET, Mono, digital television}

\maketitle

\ifincludesimulation
\begin{center}
\fcolorbox{red}{yellow!15}{\parbox{0.92\linewidth}{\centering\bfseries
SIMULATION-ONLY INTERNAL DRAFT---NOT FOR SUBMISSION. All Results values are
generated from prospective targets and must be replaced by observed runs.}}
\end{center}
\fi

%% =====================================================================
\section{Introduction}
\label{sec:intro}

Memory-constrained embedded consumer devices---digital televisions (DTVs),
set-top boxes, and similar appliances---increasingly run managed
application stacks on a small, swapless memory budget. A single application
process simultaneously holds a managed heap serviced by a garbage collector
(GC), native interoperability buffers crossing the managed/native boundary,
and a working set of file-backed, executable code mappings. Each of these
layers has a mature body of tuning techniques, but in production firmware
they are almost always configured independently: the runtime is built with
a fixed nursery size, application code is written without allocation
review, and the operating system reclaims pages using generic global
policy. Independent tuning is locally reasonable and globally fragile.
Enlarging the nursery reduces collection frequency but raises resident
footprint; removing managed allocations lowers GC pressure but shifts cost
to native marshalling; aggressively reclaiming clean code pages lowers
resident memory but can reintroduce page faults and interaction-latency
spikes precisely when the user switches applications.

The systems question we address is therefore \emph{not} whether larger
nursery sizes, value types, \texttt{madvise}-based reclamation, or
lightweight online predictors work in isolation. Each is established. The
question is whether these mechanisms can be \emph{coordinated} under a
single embedded-process memory budget without trading fewer collections
for interop allocation, or lower resident memory for refault-induced
interaction latency. This framing makes SAGE-MM a measured integration and
control study, not a claim that its constituent mechanisms are novel.

We deliberately separate the three interventions by \emph{lifecycle} and
\emph{adaptivity} (Section~\ref{sec:boundary}). Heap sizing is a
runtime-build decision; value-type conversion is a reviewed
source-and-recompile decision; only reclamation scheduling and compaction
gating change online. Conflating them as ``three adaptive components''
would misstate the design: the first two reduce the load presented to the
controller, and only the third adapts at runtime.

\paragraph{Contributions.}
This paper makes four bounded contributions.
\begin{enumerate}
  \item It characterizes managed-heap, interop-allocation, and file-backed
  executable-page pressure under identical application transitions on the
  evaluated firmware, and reports architecture-specific rather than
  universal effects (Section~\ref{sec:motivation}).
  \item It specifies a cross-layer action boundary: heap sizing is a
  runtime-build decision, value-type conversion is a reviewed source
  decision, and only reclamation scheduling and compaction gating change
  online (Section~\ref{sec:boundary}).
  \item It provides an implementable controller definition with
  dimensionless features and target, a causal update order, bounded
  actions, hysteresis, cooldown, fail-closed telemetry handling, and an
  explicit non-learning threshold fallback (Section~\ref{sec:design}).
  \item It defines a falsifiable evaluation---a full-factorial ablation
  with held-out policy comparison, adverse-workload stress, and endurance
  runs---and reports per-mechanism cost and interaction effects
  (Sections~\ref{sec:method}--\ref{sec:results}).
\end{enumerate}

We make no claim that SAGE-MM is a new garbage collector, or that results
from two DTV models establish all embedded .NET deployments. Claims outside
the measured device, memory, runtime revision, and workload matrix are
stated as hypotheses for future evaluation.

%% =====================================================================
\section{Background and Motivation}
\label{sec:motivation}

\subsection{Embedded managed runtimes under a shared budget}
The deployment studied here is a vendor fork of the Mono runtime inside a
Tizen .NET~6 firmware image, running on DTV-class hardware: a 1\,GiB ARM32
device and a 3\,GiB ARM64 device~\cite{monoruntime}. Mono and CoreCLR are
distinct runtime
implementations; ``Mono CoreCLR'' is not a runtime name, and the studied
deployment must be identified by its vendor Mono source revision and patch
hashes rather than by a generic runtime label. The devices are swapless:
when the process working set exceeds the cgroup budget, the outcome is an
out-of-memory (OOM) kill or a watchdog restart, not paging to disk. This
makes resident footprint a hard correctness constraint, not merely a
performance metric, and makes the interaction between GC, interop, and
code-page residency directly observable at the process level.

\subsection{Why independent tuning regresses}
Three cross-layer tensions motivate coordination.
\begin{itemize}
  \item \textbf{Heap vs.\ footprint.} A larger initial managed allocation
  reduces early collections and compaction cost, most visibly on ARM32,
  but raises steady-state resident memory. The effect is
  architecture-specific and must be measured per build, not assumed.
  \item \textbf{Allocation vs.\ interop.} Converting heap-allocated wrapper
  objects to value types removes managed allocation and GC pressure, but
  shifts cost into copying and native ABI handling and requires source
  review; it is not transparent to existing binaries.
  \item \textbf{Residency vs.\ latency.} Dropping clean, file-backed code
  pages with \texttt{madvise(MADV\_DONTNEED)} lowers resident footprint but
  can force refaults on immediate hot-code reuse, degrading input and frame
  tail latency during rapid application switching.
\end{itemize}
The central hypothesis is that a guarded, online policy can hold the
reclamation benefit while bounding the refault penalty, and that the two
static interventions reduce the pressure the controller must manage. This
is falsifiable: uncoordinated static reclamation is actually measured to inflate the
fault rate, whereas guarded coordination should hold it near a bounded
target (Section~\ref{sec:results}).

%% =====================================================================
\section{System Boundary and Action Space}
\label{sec:boundary}

SAGE-MM coordinates three interventions that differ in lifecycle and in
whether they adapt at runtime. Table~\ref{tab:boundary} states the boundary
precisely; conflating these layers is the single most common source of
overclaiming in this design space.

\begin{table}[t]
  \caption{Cross-layer action boundary. Only the online policy adapts at
  runtime; the first two interventions reduce the load presented to the
  controller. ``Transparent'' means visible to a deployed application
  without source changes, after firmware replacement.}
  \label{tab:boundary}
  \small
  \begin{tabular}{@{}p{3.1cm}p{2.4cm}p{1.7cm}p{4.4cm}@{}}
    \toprule
    \textbf{Mechanism} & \textbf{Lifecycle} & \textbf{Adaptive?} &
    \textbf{Notes} \\
    \midrule
    Heap build config
    (\texttt{INITIAL\_ALLOC}, e.g.\ 16$\to$64\,MiB, ARM32) &
    Vendor runtime build & No &
    Requires rebuilding the vendor Mono runtime; transparent after firmware
    replacement; not an online control. \\
    Value-type interop conversion &
    Source / recompile & No &
    Reviewed class-to-struct migration; changes identity, boxing, layout,
    and native ABI; \emph{not} transparent to existing binaries. \\
    Reclamation \& compaction control &
    Runtime, online & \textbf{Yes} &
    The controller changes only the reclamation interval and the compaction
    gate exposed by the runtime host. \\
    \bottomrule
  \end{tabular}
\end{table}

Two consequences follow. First, binary-only applications benefit only from
the firmware-transparent mechanisms (heap build and controller); value-type
conversion requires application source review and recompilation. Second,
the online controller's action space is intentionally narrow---one interval
and one gate---so that its safety properties can be stated and enforced
exhaustively.

%% =====================================================================
\section{Design}
\label{sec:design}

We present the design in three parts, matching the action boundary:
static interventions (\S\ref{sec:static}), the adaptive controller
(\S\ref{sec:controller}), and the reclamation implementation
(\S\ref{sec:reclaim}). Figure~\ref{fig:arch} shows how the three layers
compose at runtime.

\begin{figure}[t]
  \centering
  \setlength{\fboxsep}{7pt}%
  \fbox{\parbox[c][4.6cm][c]{0.9\linewidth}{\centering\small
  \textbf{[Figure placeholder --- redraw as a labeled block diagram]}\\[5pt]
  System overview. Per-interval telemetry
  ($L_{\mathrm{gc}}$, $F_h$, $P_f$, $\Delta M$) is normalized
  (Eq.~\ref{eq:pressure}) and consumed by the online controller, which
  emits two bounded actions: the reclamation interval $T$ and the compaction
  gate. The reclamation path ranks cold, private, file-backed code mappings
  and issues \texttt{madvise(MADV\_DONTNEED)} through the conservative native
  helper, gated by the hot-reuse and cooldown guards. Static interventions
  (heap build $G$, interop $I$) sit upstream and reduce the pressure
  presented to the controller; they do not adapt online.}}
  \caption{SAGE-MM runtime composition and the narrow online action space
  (reclamation interval $T$ and compaction gate). Only the control path
  adapts; heap build and interop conversion are fixed before deployment.}
  \label{fig:arch}
\end{figure}

\subsection{Static interventions}
\label{sec:static}

\paragraph{Architecture-specific heap build.}
The managed runtime is rebuilt with an architecture-specific initial
allocation. On the memory-tight ARM32 device, raising the initial nursery
reduces early collection and compaction work; on ARM64 the trade-off point
differs. Because the setting is compiled into the vendor runtime, it is a
build-time decision recorded by patch hash and build flags, not an online
knob. Its effect on peak resident memory and collection count is measured
per architecture (RQ1) rather than assumed to transfer.

\paragraph{Value-type interop conversion.}
A source-time analyzer (diagnostic \textsf{DTV0001}) flags plain-old-data
wrapper classes whose conversion to \texttt{struct} removes heap allocation
across the managed/native boundary. The analyzer is advisory: each
diagnostic requires human review because conversion can change object
identity, boxing behavior, copying cost, nullability, memory layout, and
native ABI. It must not be described as requiring no application
modification~\cite{dotnetinterop}. The intervention reduces allocation rate
and GC pressure
(RQ2), lowering the load presented to the controller.

\subsection{Adaptive controller}
\label{sec:controller}

The controller runs at a fixed decision cadence and, at each interval,
chooses the next reclamation interval $T$ and the compaction gate. All
initialization values, feature scales, thresholds, and update constants are
externalized in a single options block so that no magic constant is buried
in code. Algorithm~\ref{alg:controller} states the complete per-interval
procedure; the paragraphs below define each step.

\paragraph{Telemetry and normalization.}
Each decision uses the preceding interval's GC pause
$L_{\mathrm{gc}}$ (ms), fragmentation ratio $F_h$, page-fault rate $P_f$
(minor plus major faults per second, divided by the \emph{measured}
wall-clock sample duration rather than an assumed one second), and positive
resident growth $\Delta M$ (MiB). Features are dimensionless and clamped to
$[0,2]$:
\begin{equation}
\mathbf{x} = \Bigl[\,1,\;
  \mathrm{clamp}\!\bigl(\tfrac{L_{\mathrm{gc}}}{30}\bigr),\;
  \mathrm{clamp}\!\bigl(\tfrac{F_h}{0.12}\bigr),\;
  \mathrm{clamp}\!\bigl(\tfrac{P_f}{100}\bigr),\;
  \mathrm{clamp}\!\bigl(\tfrac{\max(0,\Delta M)}{50}\bigr)\Bigr].
\end{equation}
The observed dimensionless pressure is the maximum normalized
service-objective violation,
\begin{equation}
\label{eq:pressure}
y = \mathrm{clamp}\!\Bigl(
  \max\bigl(\tfrac{L_{\mathrm{gc}}}{30\,\text{ms}},\;
            \tfrac{P_f}{100\,\text{s}^{-1}},\;
            \tfrac{\max(0,\Delta M)}{50\,\text{MiB}}\bigr),\,0,\,2\Bigr),
\end{equation}
so that no coefficient silently combines milliseconds, counts, and MiB, and
a breach in any single dimension is actionable.

\paragraph{Policies.}
Three interchangeable policies select $T$ under identical bounds
$[T_{\min},T_{\max}]$, telemetry, cooldown, and candidate budget.
\begin{itemize}
  \item \emph{Threshold} (non-learning): shorten $T$ by 20\% when
  $y>1.1$, lengthen it by 20\% when $y<0.9$, otherwise hold. This is a
  \emph{tuned} comparator, not an untuned default.
  \item \emph{EWMA}: $T' = \beta T + (1-\beta)\,T/\mathrm{clamp}(y,0.5,2)$
  with $\beta = 0.85$.
  \item \emph{Ridge} (online): predict next-interval pressure
  $\hat{y}(t{+}1) = \mathrm{clamp}(\mathbf{w}(t)^{\!\top}\mathbf{x}(t),0,2)$
  and set $T' = T/\mathrm{clamp}(0.5+\hat{y},0.5,1.5)$, bounded to
  $[T_{\min},T_{\max}]$.
\end{itemize}

\paragraph{Causal learning.}
The ridge prediction is computed \emph{before} its update. When the true
$y(t{+}1)$ arrives, the controller records the causal prequential loss
$[\hat{y}(t{+}1)-y(t{+}1)]^2/2$ and performs one online ridge SGD step with
the pending feature vector:
\begin{equation}
w_i(t{+}1) \leftarrow w_i(t) -
  \eta\,\bigl[x_i(t)\,(\hat{y}(t{+}1)-y(t{+}1)) + \lambda\, w_i(t)\bigr],
\end{equation}
where the bias weight starts at $1$ and the others at $0$,
$\eta = 5\times10^{-4}$, and $\lambda = 10^{-4}$. Training occurs only on a
designated tuning trace; reporting calls pass \texttt{updateModel:false} to
freeze weights, and \texttt{BeginReporting()} clears the training-boundary
pending prediction while retaining learned weights. Training and reporting
traces are disjoint.

\paragraph{Compaction gating and hysteresis.}
In adaptive modes, compaction is disabled below fragmentation $0.05$ and
forcibly re-enabled at $0.12$ or after three deferrals. This two-threshold
guard prevents compaction starvation. Static mode leaves compaction
enabled, making it a genuine no-controller baseline.

\paragraph{Fail-closed robustness.}
Every sample is validated before feature processing. Non-finite values,
negative pauses or fault rates, and fragmentation outside $[0,1]$ trigger a
fail-closed decision: hold the previous bounded interval, keep compaction
enabled, suppress reclamation, record \textsf{InvalidTelemetry}, and skip
the model update. Persistent prediction clipping at $0$ or $2$ reports
\textsf{ModelSaturation}, clears the pending ML update, and uses the
predeclared threshold action. Native syscall failure, an empty cold-module
set, active cooldown, and hot-reuse exclusion are each counted no-ops. The
evaluation reports the frequency, duration, and outcome of every fallback
class.

\begin{algorithm}[t]
\caption{SAGE-MM controller decision for one interval. Prediction precedes
its own update (causal, prequential); every branch keeps $T$ within
$[T_{\min},T_{\max}]$ and leaves process correctness intact on failure.}
\label{alg:controller}
\small
\begin{algorithmic}[1]
\Require raw sample $s=(L_{\mathrm{gc}},F_h,P_f,\Delta M,\mathrm{dt})$; state
$(T,\mathbf{w},\textit{mode})$; bounds $[T_{\min},T_{\max}]$
\Ensure next reclamation interval $T$ and compaction gate
\If{$\neg\textsc{Valid}(s)$}\Comment{non-finite, negative rate/pause, or $F_h\notin[0,1]$}
  \State record \textsf{InvalidTelemetry}; enable compaction; suppress reclamation
  \State \Return $\mathrm{clamp}(T,T_{\min},T_{\max})$ \Comment{fail-closed: hold, no model update}
\EndIf
\State $\mathbf{x}\gets\textsc{Normalize}(s)$; \quad $y\gets\textsc{Pressure}(s)$ \Comment{Eq.~(1),~(2)}
\If{$\textit{mode}=\textsf{Ridge}$}
  \State $\hat{y}\gets\mathrm{clamp}(\mathbf{w}^{\!\top}\mathbf{x},0,2)$
  \If{$\hat{y}$ persistently clipped at $0$ or $2$}
    \State record \textsf{ModelSaturation}; clear pending update; apply \textsf{Threshold} action
  \Else
    \State $T\gets T/\mathrm{clamp}(0.5+\hat{y},0.5,1.5)$
  \EndIf
\ElsIf{$\textit{mode}=\textsf{EWMA}$}
  \State $T\gets \beta T+(1-\beta)\,T/\mathrm{clamp}(y,0.5,2)$
\ElsIf{$\textit{mode}=\textsf{Threshold}$}
  \State $T\gets T\cdot(1-0.2)$ if $y>1.1$; $\;T\cdot(1+0.2)$ if $y<0.9$; else hold
\EndIf
\State $T\gets\mathrm{clamp}(T,T_{\min},T_{\max})$
\State \textsc{UpdateGate}$(F_h)$ \Comment{off below $0.05$; force on at $0.12$ or after 3 deferrals}
\If{$\neg\textsc{InCooldown}$ \textbf{and} reclamation enabled}
  \State rank modules by $\mathrm{Cold}(\cdot)$; take smallest prefix $K$ meeting the byte budget
  \State \textsc{FlushModules}$(K)$ excluding recently accessed (hot-reuse guard)
\EndIf
\If{$\textit{mode}=\textsf{Ridge}$ \textbf{and} a pending prediction exists \textbf{and} on tuning trace}
  \State record prequential loss $[\hat{y}-y]^2/2$; take one ridge SGD step \Comment{Eq.~(4)}
\EndIf
\State \Return $T$
\end{algorithmic}
\end{algorithm}

\subsection{Reclamation implementation}
\label{sec:reclaim}

\paragraph{Coldness score and budget-driven selection.}
For an eligible module $a$ that has been idle for at least a configured
guard period, the coldness score combines recency, frequency, and size:
\begin{equation}
\mathrm{Cold}(a) = 0.6\,\frac{\mathrm{age}(a)}{\mathrm{maxAge}}
  + 0.3\,\Bigl[1-\frac{\mathrm{accesses}(a)}{\mathrm{maxAccesses}}\Bigr]
  + 0.1\,\frac{\mathrm{cleanBytes}(a)}{\mathrm{maxCleanBytes}}.
\end{equation}
Candidates are ranked in descending score; the number reclaimed, $K$, is
the smallest prefix whose cumulative clean bytes reach a configured byte
budget---it is not fixed at five. A recent-access exclusion is the
hot-reuse safety guard.

\paragraph{Conservative native helper.}
The native helper considers page-aligned ranges from
\texttt{/proc/self/maps}, accepts private \texttt{r-{}-p}/\texttt{r-xp}
file mappings, and rejects anonymous, writable, shared, and deleted-file
mappings. \texttt{madvise(MADV\_DONTNEED)} never unmaps a range: a later
access can fault file content back in. Each syscall result is checked; zero
means no candidate, a positive value is the number dropped, and a negative
value is $-\text{errno}$~\cite{madvise}. Callers retain the mapping after
failure. Because
\texttt{maps} lacks per-page residency and dirty state, a production port
must additionally verify \texttt{Private\_Clean} from
\texttt{/proc/self/smaps}, use \texttt{mincore} for residency, serialize
against module unload, and allowlist executable modules.

%% =====================================================================
\section{Implementation and Reproducibility Boundary}
\label{sec:impl}

We release a reference reproducibility kit that implements the controller,
telemetry collectors, the policy enforcer, the conservative native helper,
and the advisory analyzer, with a demo workload that simulates application
switches and allocation bursts. The kit is a user-space analogue: its
forced Gen-0 pause, fragmentation estimate, and file-length size proxy are
\emph{not} the production vendor patch, production access-event stream, or
\texttt{Private\_Clean} attribution, and must not be used as evidence for
the paper's memory, latency, or reliability claims. The commercial target
was a vendor Mono revision inside Tizen .NET~6 firmware; a publishable
artifact records the vendor source commit, patch hashes, compiler,
kernel/firmware identifiers, ARM32/ARM64 build flags, and the exact hosting
interfaces used. Separating the reproducible control logic from the
proprietary firmware integration is deliberate: it lets the control-flow,
safety guards, and statistics be validated openly without claiming binary
transparency for a stock runtime.

%% =====================================================================
\section{Evaluation Methodology}
\label{sec:method}

This section specifies the measurement protocol. All manuscript numbers are
regenerated from one versioned observed-result bundle; no number is inferred
by subtracting results collected under different device states.

\subsection{Research questions}
\begin{description}
  \item[RQ1] Does architecture-specific heap configuration change GC and
  compaction cost, and by how much per architecture?
  \item[RQ2] Do value-type interop conversion and guarded reclamation
  reduce footprint safely (without unacceptable allocation shift or
  refault)?
  \item[RQ3] Does online coordination improve robustness (footprint and
  tail latency) relative to fixed, threshold, and EWMA policies on held-out
  traces?
\end{description}
Table~\ref{tab:rqmap} binds each research question to the experiment that
answers it and to the primary outcomes reported for it, so that every result
subsection traces back to a stated question and no reported number is
orphaned from a hypothesis.

\begin{table}[t]
  \caption{Research-question traceability. Each RQ maps to a specific
  experiment and a fixed set of primary outcomes; Section~\ref{sec:results}
  reports results under the same grouping.}
  \label{tab:rqmap}
  \small
  \begin{tabular}{@{}p{0.9cm}p{4.0cm}p{5.6cm}@{}}
    \toprule
    \textbf{RQ} & \textbf{Experiment} & \textbf{Primary outcomes} \\
    \midrule
    RQ1 & ARM32/ARM64 heap-configuration sweep &
    Peak/mean PSS, collection count, p50/p95/p99/max GC pause, compaction
    work, per architecture. \\
    RQ2 & Identical interop micro/macro traces with and without value-type
    conversion &
    Allocated bytes, allocation rate, PSS, throughput, native ABI
    correctness checks. \\
    RQ3 & Full-factorial ablation plus hot-reuse, rapid-switch, and
    slow-storage stress on held-out traces &
    Reclaimed \texttt{Private\_Clean} bytes, minor/major faults, reload and
    storage latency, frame/input p95/p99, controller CPU, with run-level
    95\% CIs. \\
    \bottomrule
  \end{tabular}
\end{table}

\subsection{Platforms and workloads}
The main evaluation uses the 1\,GiB ARM32 and 3\,GiB ARM64 DTV-class
devices plus at least one independently assembled constrained embedded
Linux platform (for example, an ARM64 single-board computer with a
restricted cgroup memory budget). For every platform we report SoC and core
count, ISA, physical and cgroup memory, swap configuration, storage model,
kernel, runtime name and commit, GC patch hash, firmware image, compiler,
and thermal/power policy. Workloads span steady playback, allocation burst,
interop-intensive, hot-code reuse, rapid switching (1/5/10\,s), slow-storage
injection, and a multi-hour endurance script, each with published state
transitions and seeds.

\subsection{Baselines}
We use canonical names without aliases in every table, figure, and
paragraph: \textsf{Stock}, \textsf{Static-G}, \textsf{Static-GI},
\textsf{Static-GIR}, \textsf{Threshold-GIR}, \textsf{EWMA-GIR}, and
\textsf{Ridge-GIR}. Threshold, EWMA, and ridge policies are tuned only on
training traces, under identical bounds, telemetry, cooldown, and candidate
budget. Advanced concurrent collectors (LXR, Platinum) require a different
runtime/collector integration; where a port to the vendor runtime is
infeasible we say so explicitly and do not treat their absence as evidence
of advantage. Storage GC systems (AGC, DumpKV) are conceptual
adaptive-policy precedents, not executable managed-heap baselines
(Section~\ref{sec:related}).

\subsection{Full-factorial ablation}
\label{sec:method-factorial}
We run all 16 combinations of heap build ($G$), interop ($I$), reclamation
($R$), and adaptive controller ($C$), from the default \texttt{0000} to the
full \texttt{1111}, under identical workloads. For every row we report
independent runs $n$, duration, peak/mean PSS, allocation rate, GC count
and p50/p95/p99/max pause, reclaimed private-clean bytes, minor/major
faults, reload and storage latency, frame/input p95/p99, controller CPU and
allocation, and syscall duration/failures, each with a 95\% bootstrap
confidence interval. We analyze main and interaction effects; actual
impacts are never labeled as ablation results.

\subsection{Adverse and endurance protocol}
Reclamation adversity tests cover cold steady state, immediate hot-page
reuse, 1/5/10\,s switching, slow-storage injection, concurrent execution of
targeted mappings, \texttt{madvise} failure injection, and a fault storm
above 500\,faults/s. We collect \texttt{Private\_Clean} before/after from
\texttt{smaps}, faults, bytes read, reload latency, syscall errno/duration,
frame drops, input latency, and guard activations, and verify that no
writable, anonymous, or shared mapping is selected and that failure leaves
process correctness intact. Endurance evidence requires multiple
independent eight-hour runs per primary platform with OOM/kernel logs,
censoring rules, and per-run trajectories; five short runs do not support a
production-stability claim.

\subsection{Statistics}
We use at least the preregistered number of independent, reset/seeded runs
per condition (a documented minimum of 30 for short workloads unless a power
analysis supports another $n$), construct two-sided run-level percentile
bootstrap confidence intervals (10{,}000 resamples), and never treat
within-run samples as independent. Treatment order is randomized;
firmware, thermal, and network state are locked; training and reporting
traces are disjoint.

%% =====================================================================
\section{Results}
\label{sec:results}

\paragraph{Reporting discipline.}
The actual table below is followed by a fixed-seed simulation-only dry run
that validates the planned analysis and page layout. No value in the dry run
is an observation or achieved effect. The submission version replaces the
entire generated fragment from one versioned observed-result bundle and names
each baseline, workload aggregation, statistic, denominator, and confidence
interval; RSS, PSS, managed heap, and \texttt{Private\_Clean} are never
interchanged.

\ifincludesimulation
\subsection{Preregistered actual measurements (hypotheses, not results)}
\label{sec:actual}
To keep the design falsifiable, Table~\ref{tab:actual} states
preregistered engineering \emph{actual measurements} normalized to
$\textsf{Stock}=100$ (lower is better). These are hypotheses fixed before
the additional experiments; they are \textbf{not} measured outcomes and are
never quoted in the abstract or conclusion. The final paper adds adjacent
\textsf{Observed effect}, \textsf{95\% CI}, $n$, and \textsf{Pass/Fail}
columns generated from the observed bundle.

\begin{table}[t]
  \caption{\textbf{Preregistered actual measurements, not measured results.}
  Normalized to $\textsf{Stock}=100$; lower is better. The key falsifiable
  actual measurement is that unguarded static reclamation (\textsf{Static-GIR})
  inflates the fault index to $\approx$160, while guarded policies hold it
  near 122--132. If a policy's 95\% CI overlaps \textsf{EWMA-GIR} or its
  controller CPU exceeds 1.5\%, the paper must report that ML offers no
  demonstrated advantage.}
  \label{tab:actual}
  \small
  \begin{tabular}{@{}lrrrrr@{}}
    \toprule
    \textbf{Baseline} & \textbf{Peak PSS} & \textbf{GC p99} &
    \textbf{Fault} & \textbf{Input p99} & \textbf{Ctrl.\ CPU} \\
    \midrule
    \textsf{Stock}         & 100 & 100 & 100 & 100 & 0.0\% \\
    \textsf{Static-G}      &  94 &  75 & 100 &  90 & $\le$0.1\% \\
    \textsf{Static-GI}     &  87 &  71 & 100 &  86 & $\le$0.2\% \\
    \textsf{Static-GIR}    &  79 &  71 & 160 &  91 & $\le$0.4\% \\
    \textsf{Threshold-GIR} &  78 &  66 & 132 &  85 & $\le$0.8\% \\
    \textsf{EWMA-GIR}      &  77 &  64 & 125 &  82 & $\le$1.0\% \\
    \textsf{Ridge-GIR}     &  76 &  62 & 122 &  80 & $\le$1.5\% \\
    \bottomrule
  \end{tabular}
\end{table}

\subsection{Simulation-only pipeline validation}
\label{sec:simulated-results}
The following generated fragment fills every planned RQ, factorial,
additional-platform, adverse, and endurance result field with fixed-seed
\emph{synthetic} samples. These values are not observations and cannot be
used in the abstract, conclusions, or submission. Their sole purpose is to
exercise the analysis and page-layout pipeline before device measurements.
% AUTO-GENERATED by scripts/generate_simulated_results.py
% SYNTHETIC PIPELINE VALIDATION ONLY -- NOT OBSERVED DATA.
\subsection{Simulation-only RQ1 dry run: architecture-specific heap configuration}
\textbf{Synthetic data---not observed.} The seeded model projects the heap build to reduce normalized GC p99 from 100 to 75 and instantiates the 2.6$\times$ ARM32/ARM64 compaction-frequency hypothesis solely to validate the planned figure and prose.
\begin{figure}[t]
  \centering\small
  \begin{tabular}{@{}lr@{ }l@{}}
  Stock & 100 & \textcolor{blue}{\rule{50mm}{2.4mm}}\\
  Static-G & 75 & \textcolor{blue}{\rule{37.5mm}{2.4mm}}\\
  \end{tabular}
  \caption{\textbf{SIMULATED---NOT OBSERVED.} RQ1 GC-p99 layout dry run, normalized to Stock=100. Synthetic values cannot validate the runtime.}
  \Description{Synthetic horizontal bars compare normalized GC-p99 indices 100 for Stock and 75 for Static-G.}
  \label{fig:sim-rq1}
\end{figure}

\subsection{Simulation-only RQ2 dry run: representation and reclamation}
\textbf{Synthetic data---not observed.} The model projects interop allocation index 100$\to$72 and static-reclamation PSS index 100$\to$92, while deliberately raising its fault index to 160. This encodes the refault risk rather than concealing it.
\begin{figure}[t]
  \centering\small
  \begin{tabular}{@{}lr@{ }l@{}}
  Allocation, Stock & 100 & \textcolor{blue}{\rule{31mm}{2.4mm}}\\
  Allocation, Interop & 72 & \textcolor{blue}{\rule{22.3mm}{2.4mm}}\\
  Faults, Stock & 100 & \textcolor{red}{\rule{31mm}{2.4mm}}\\
  Faults, static reclaim & 160 & \textcolor{red}{\rule{49.6mm}{2.4mm}}\\
  \end{tabular}
  \caption{\textbf{SIMULATED---NOT OBSERVED.} RQ2 allocation/refault trade-off layout.}
  \Description{Synthetic bars show allocation index falling from 100 to 72 and fault index rising from 100 to 160 under static reclamation.}
  \label{fig:sim-rq2}
\end{figure}

\subsection{Simulation-only RQ3 dry run: coordinated policy}
\textbf{Synthetic data---not observed.} The mock input-p99 indices are Static-GIR 91, Threshold-GIR 85, EWMA-GIR 82, and Ridge-GIR 80. The small Ridge--EWMA gap predeclares that overlapping real confidence intervals imply no demonstrated ML advantage.
\begin{figure}[t]
  \centering\small
  \begin{tabular}{@{}lr@{ }l@{}}
  Static-GIR & 91 & \textcolor{blue}{\rule{45.5mm}{2.4mm}}\\
  Threshold-GIR & 85 & \textcolor{blue}{\rule{42.5mm}{2.4mm}}\\
  EWMA-GIR & 82 & \textcolor{blue}{\rule{41mm}{2.4mm}}\\
  Ridge-GIR & 80 & \textcolor{blue}{\rule{40mm}{2.4mm}}\\
  \end{tabular}
  \caption{\textbf{SIMULATED---NOT OBSERVED.} RQ3 controller-comparison layout.}
  \Description{Synthetic bars compare input-p99 indices 91, 85, 82, and 80 for Static-GIR, Threshold-GIR, EWMA-GIR, and Ridge-GIR.}
  \label{fig:sim-rq3}
\end{figure}

\begin{table*}[t]
\caption{\textbf{SIMULATED---NOT OBSERVED.} Full $2^4$ pipeline dry run ($n=30$ generated samples/cell; 10,000 seeded bootstrap resamples). E=prospective actual index; S=simulated mean [95\% CI]. Synthetic PASS tests only the reporting pipeline.}
\label{tab:sim-factorial}
\centering\scriptsize
\resizebox{\textwidth}{!}{%
\begin{tabular}{@{}lrrrrrrrrr@{}}
\toprule
Cell & PSS E & PSS S [CI] & GC E & GC S [CI] & Fault E & Fault S [CI] & Input E & Input S [CI] & Status\\
\midrule
G0I0R0C0 & 100.0 & 100.1 [99.6, 100.7] & 100.0 & 100.2 [99.3, 101.1] & 100.0 & 99.1 [96.4, 101.7] & 100.0 & 100.1 [99.2, 101.0] & PASS\\
G0I0R0C1 & 98.0 & 98.6 [97.8, 99.3] & 90.0 & 90.2 [89.4, 91.0] & 100.0 & 101.1 [98.6, 103.4] & 92.0 & 92.2 [91.6, 92.8] & PASS\\
G0I0R1C0 & 92.0 & 91.5 [90.8, 92.2] & 100.0 & 100.5 [99.5, 101.5] & 160.0 & 159.7 [157.8, 161.7] & 105.0 & 105.2 [104.3, 106.0] & PASS\\
G0I0R1C1 & 90.0 & 89.6 [88.9, 90.3] & 90.0 & 90.4 [89.4, 91.3] & 125.0 & 125.3 [122.7, 127.9] & 96.0 & 96.6 [95.8, 97.5] & PASS\\
G0I1R0C0 & 93.0 & 92.4 [91.7, 93.2] & 95.0 & 94.8 [94.0, 95.6] & 100.0 & 100.1 [97.9, 102.3] & 96.0 & 96.4 [95.5, 97.3] & PASS\\
G0I1R0C1 & 91.0 & 90.6 [90.1, 91.2] & 86.0 & 86.3 [85.5, 87.2] & 100.0 & 100.0 [97.9, 102.0] & 89.0 & 88.6 [87.7, 89.5] & PASS\\
G0I1R1C0 & 85.0 & 85.0 [84.1, 85.9] & 95.0 & 95.8 [95.0, 96.5] & 160.0 & 158.5 [156.5, 160.6] & 101.0 & 101.3 [100.6, 102.1] & PASS\\
G0I1R1C1 & 83.0 & 82.9 [82.4, 83.5] & 86.0 & 85.7 [84.7, 86.6] & 125.0 & 126.2 [123.7, 128.8] & 92.0 & 91.6 [90.9, 92.3] & PASS\\
G1I0R0C0 & 94.0 & 94.5 [93.7, 95.3] & 75.0 & 74.5 [73.7, 75.3] & 100.0 & 103.6 [101.2, 105.9] & 90.0 & 90.2 [89.3, 91.1] & PASS\\
G1I0R0C1 & 92.0 & 91.7 [90.9, 92.6] & 68.0 & 67.8 [66.9, 68.7] & 100.0 & 97.4 [94.8, 100.1] & 83.0 & 82.8 [82.1, 83.5] & PASS\\
G1I0R1C0 & 86.0 & 85.9 [85.1, 86.6] & 75.0 & 74.5 [73.5, 75.5] & 160.0 & 162.3 [160.0, 164.7] & 95.0 & 93.8 [92.8, 94.7] & PASS\\
G1I0R1C1 & 84.0 & 84.2 [83.4, 85.0] & 68.0 & 68.4 [67.6, 69.3] & 125.0 & 126.9 [124.4, 129.3] & 86.0 & 85.9 [85.0, 86.7] & PASS\\
G1I1R0C0 & 87.0 & 86.5 [85.9, 87.2] & 71.0 & 71.2 [70.3, 72.0] & 100.0 & 99.3 [97.0, 101.6] & 86.0 & 86.8 [86.0, 87.5] & PASS\\
G1I1R0C1 & 85.0 & 84.6 [83.8, 85.5] & 64.0 & 64.0 [63.0, 64.9] & 100.0 & 100.4 [97.9, 102.8] & 79.0 & 79.9 [79.2, 80.6] & PASS\\
G1I1R1C0 & 79.0 & 79.3 [78.6, 80.1] & 71.0 & 71.1 [70.5, 71.7] & 160.0 & 160.3 [157.9, 162.6] & 91.0 & 90.8 [90.0, 91.6] & PASS\\
G1I1R1C1 & 77.0 & 76.8 [76.1, 77.6] & 64.0 & 64.0 [63.1, 65.0] & 125.0 & 126.3 [124.0, 128.6] & 82.0 & 81.4 [80.6, 82.2] & PASS\\
\bottomrule
\end{tabular}}
\end{table*}

\subsection{Simulation-only additional-platform dry run}
\begin{table}[t]
\caption{\textbf{SIMULATED---NOT OBSERVED.} Additional-platform layout.}
\label{tab:sim-platform}
\centering\scriptsize
\begin{tabular}{@{}lrrrrr@{}}\toprule
Platform & PSS & GC p99 & Fault & Input p99 & CPU\\\midrule
Constrained ARM64 SBC (synthetic) & 80.2 & 68.4 & 129.1 & 84.6 & 1.3\%\\
\bottomrule\end{tabular}
\end{table}

\subsection{Simulation-only adverse and endurance dry run}
\begin{table}[t]
\caption{\textbf{SIMULATED---NOT OBSERVED.} Stress/endurance layout.}
\label{tab:sim-adverse}
\centering\scriptsize
\begin{tabular}{@{}p{3.0cm}p{3.2cm}p{2.2cm}l@{}}\toprule
Scenario & Synthetic outcome & Threshold & Status\\\midrule
Hot reuse, normal & reload p99 22.1 ms & $\leq25$ ms & PASS\\
Hot reuse, slow storage & reload p99 69.3 ms & $\leq75$ ms & PASS\\
Rapid switching & fault index 131.0 & $\leq135$ & PASS\\
Failure injection & 0 correctness failures & 0 & PASS\\
Fault storm & 100\% reclaim suppressed & 100\% & PASS\\
Endurance & 10$\times$8 h; 0 mock OOM/watchdog & 80 h; 0 OOM & PASS\\
\bottomrule\end{tabular}
\end{table}

\paragraph{Mandatory replacement.}
All values above are sampled from a seeded probability model centered on the actual measurements. They cannot expose firmware, runtime, measurement, safety, or performance defects. This entire generated fragment must be replaced by provenance-backed runs before submission.


\else
%% This file is deliberately absent until the observed result bundle has been
%% summarized and reviewed. A submission build fails rather than silently
%% falling back to mock data.
% AUTO-GENERATED by code/scripts/make_observed_results_tex.py
% OBSERVED PROXY-HARNESS DATA -- NOT VENDOR DTV Mono/.NET6 FIRMWARE EVIDENCE.
\subsection{Observed proxy-harness factorial (ARM64 Linux; not firmware evidence)}
\label{sec:observed-results}
\noindent\textbf{Observed, non-synthetic --- but a proxy harness, not firmware.} The values below are actually measured (\texttt{data\_status=observed}, \texttt{synthetic=false}), but they were collected by a Python \texttt{resource}~$+$~\texttt{/proc/self/smaps\_rollup}~$+$~\texttt{monotonic\_ns} harness on ARM64 Linux (kernel~6.18.35), not on a vendor DTV Mono/.NET6 runtime. The supplied provenance is \emph{partial}: the firmware vendor, version, and date are \texttt{unavailable} and are not fabricated here. GC~p99 is a Python \texttt{gc.collect} proxy, not a managed-runtime pause; PSS is the process-wide \texttt{smaps\_rollup} share, never per-mapping \texttt{Private\_Clean} reclaimed bytes. Each cell has 3 independent runs (not the preregistered $\ge$30), and per-run wall time is a sub-10\,ms micro-run. These numbers therefore validate the analysis pipeline on real data; they are \textbf{not} an achieved memory, latency, or reliability effect and are never quoted in the abstract or conclusion.

\begin{table}[t]
  \caption{\textbf{Observed proxy-harness factorial --- not firmware evidence.} Run-level means with fixed-seed ($20260829$) 10{,}000-resample 95\% bootstrap CIs over $n=3$ independent runs per cell, on the ARM64 Linux Python proxy harness. Cells are the $2^4$ \textsf{GIRC} treatments (heap, interop, reclamation, controller). Results are directionally mixed and several metrics sit near the micro-run noise floor; this table exercises the pipeline and does not substitute for the vendor-firmware campaign of Section~\ref{sec:actual}.}
  \label{tab:observed}
  \footnotesize
  \setlength{\tabcolsep}{4pt}
  \begin{tabular}{@{}lrrrr@{}}
    \toprule
    \textbf{Cell} & \textbf{Peak PSS (MiB)} & \textbf{Alloc (MiB/s)} & \textbf{GC-proxy p99 (ms)} & \textbf{Ctrl.\ CPU (\%)} \\
    \midrule
    \textsf{G0I0R0C0\,$^{\dagger}$} & 19.97~[19.97, 19.98] & 1062.3~[1028.2, 1112.5] & 0.807~[0.788, 0.832] & 0.000 \\
    \textsf{G0I0R0C1} & 12.68~[12.66, 12.71] & 1378.6~[1365.7, 1387.5] & 0.791~[0.773, 0.804] & 0.694 \\
    \textsf{G0I0R1C0} & 12.22~[12.22, 12.22] & 1396.2~[1360.1, 1414.3] & 0.810~[0.797, 0.825] & 0.000 \\
    \textsf{G0I0R1C1} & 12.22~[12.22, 12.22] & 1388.8~[1379.5, 1401.1] & 0.788~[0.785, 0.792] & 0.029 \\
    \textsf{G0I1R0C0} & 21.64~[21.60, 21.68] & 1694.1~[1653.9, 1734.2] & 0.817~[0.811, 0.827] & 0.000 \\
    \textsf{G0I1R0C1} & 18.71~[18.71, 18.71] & 1807.3~[1763.4, 1851.5] & 0.862~[0.823, 0.931] & 0.314 \\
    \textsf{G0I1R1C0} & 18.71~[18.71, 18.71] & 1903.4~[1882.7, 1941.2] & 0.860~[0.832, 0.889] & 0.000 \\
    \textsf{G0I1R1C1} & 18.72~[18.72, 18.72] & 1879.3~[1861.7, 1895.0] & 0.830~[0.811, 0.856] & 0.014 \\
    \textsf{G1I0R0C0} & 18.72~[18.72, 18.72] & 750.2~[663.2, 797.8] & 1.043~[0.829, 1.374] & 0.000 \\
    \textsf{G1I0R0C1} & 18.72~[18.72, 18.72] & 940.5~[906.0, 981.1] & 0.815~[0.769, 0.840] & 0.519 \\
    \textsf{G1I0R1C0} & 18.73~[18.72, 18.73] & 958.8~[924.7, 976.0] & 0.806~[0.773, 0.866] & 0.000 \\
    \textsf{G1I0R1C1} & 18.73~[18.73, 18.73] & 922.6~[893.5, 955.4] & 0.828~[0.782, 0.891] & 0.028 \\
    \textsf{G1I1R0C0} & 19.16~[19.11, 19.22] & 1557.6~[1525.4, 1582.1] & 0.840~[0.799, 0.898] & 0.000 \\
    \textsf{G1I1R0C1} & 18.74~[18.74, 18.74] & 1552.6~[1518.8, 1575.0] & 0.805~[0.792, 0.815] & 0.356 \\
    \textsf{G1I1R1C0} & 18.74~[18.74, 18.74] & 1546.4~[1509.8, 1566.0] & 0.841~[0.809, 0.869] & 0.000 \\
    \textsf{G1I1R1C1} & 18.74~[18.74, 18.74] & 1585.5~[1501.6, 1643.3] & 0.853~[0.810, 0.930] & 0.017 \\
    \bottomrule
  \end{tabular}\\[2pt]
  {\footnotesize $^{\dagger}$\textsf{G0I0R0C0} is the Stock reference cell. Device id prefix \texttt{bb2498ecbe619967}; collector sha256 \texttt{413a1d7e0a574cfc\ldots}}
\end{table}

\noindent Because the substrate is a proxy harness and $n=3$, the observed cells do not meet the measured-result bar in the reproducibility kit's \texttt{MEASURED\_DEVICE\_RUN\_VALUES.md} (vendor GC/EventPipe trace, per-mapping \texttt{Private\_Clean} accounting, $\ge$30 cold-reset runs, $\ge$30\,min and 8\,h durations). The vendor DTV Mono/.NET6 result cells accordingly remain \texttt{NA}; this fragment is provenance-backed observed evidence for the harness alone.


\fi

%% =====================================================================
\section{Related Work}
\label{sec:related}

We organize related work by mechanism and close each comparison with a
limitation statement rather than a generic novelty claim. Static heap
sizing, class-to-struct conversion, \texttt{madvise}, EWMA, and ridge
regression are not individually novel; the contribution is their
safety-aware coordination at the embedded Mono/Linux boundary, supported by
factorial interaction effects.

\paragraph{Managed-runtime collection and adaptive GC.}
Advanced collectors such as LXR~\cite{lxr2022} and
Platinum~\cite{platinum2020} change collection itself---concurrency,
pause/throughput objectives, and runtime integration. SAGE-MM does not
introduce a collector; it coordinates an existing vendor collector with
process-level actions (interval and gate). Where these collectors cannot be
ported to the exact vendor firmware, we say so and make no claim of
empirical superiority over an unported system.

\paragraph{Heuristic and learned controllers.}
We compare fixed intervals, a tuned pressure threshold, EWMA, and online
ridge under the same trace split and action bounds. Storage-GC learning
systems---AGC~\cite{agc2025} for burst-aware SSD flash GC and
DumpKV~\cite{dumpkv2024} for learned lifetime control in LSM key/value
separation---are conceptual precedents for lightweight learning, not
executable managed-runtime baselines; we report this domain mismatch rather
than implying empirical equivalence.

\paragraph{Application-directed reclamation.}
We compare \texttt{madvise}/discard mechanisms by candidate identification,
dirty-page safety, concurrency with unload and execution, refault behavior,
and feedback signals. SAGE-MM's differentiator is the coldness-ranked,
budget-driven selection with a hot-reuse exclusion and fail-closed guards.

\paragraph{Managed/native interop.}
We distinguish ABI/source transformations (such as value-type conversion,
which requires ownership review and recompilation) from transparent runtime
optimization, and state layout, identity, and recompilation requirements
explicitly.

%% =====================================================================
\section{Threats to Validity}
\label{sec:threats}

\paragraph{Construct validity.}
Telemetry signals are proxies for the quantities of interest. The public
kit's forced Gen-0 pause, fragmentation estimate, and file-length size proxy
are analogues, not the production GC pause, live fragmentation, or
\texttt{Private\_Clean} byte count; production evidence uses EventPipe or
vendor GC telemetry and \texttt{smaps}-verified clean bytes. We keep RSS,
PSS, managed heap, and \texttt{Private\_Clean} distinct throughout and report
the numerator, denominator, statistic, and CI for every effect.

\paragraph{Internal validity.}
Confounding across device state is the main risk. Every configuration runs
from the same rebooted image, locked thermal/power/network state, and seeded
workload script, with randomized treatment order and disjoint training and
reporting traces, so learned-policy results are never reported on their own
tuning data. No mechanism's contribution is computed by subtracting runs
collected under different states; all attributions come from the single
factorial design.

\paragraph{Conclusion validity.}
Bootstrap CIs are constructed across independent runs, never across
within-run telemetry samples, which would understate variance. We use at
least the preregistered $n$ per condition and report per-run aggregates,
OOM/kernel events, censoring, and missing-data handling, so a null or
negative result (for example, ridge not beating EWMA) is reported rather
than hidden.

\paragraph{External validity.}
Results are specific to the evaluated runtime commits, devices, memory
limits, storage configurations, and scripted workloads. The additional
independently assembled platform probes generality but does not establish it
for all embedded managed runtimes; broader applicability is stated as future
validation, not claimed.

%% =====================================================================
\section{Limitations}
\label{sec:limitations}

SAGE-MM combines known mechanisms, so its novelty is limited to their
specified coordination and empirical interaction analysis. Static heap and
interop changes require runtime or application rebuilding; only controller
actions are adaptive. \texttt{MADV\_DONTNEED} can trade PSS for page faults
and storage-dependent latency; the guards reduce but do not eliminate that
risk. The ridge model is deliberately lightweight and may offer no benefit
over a tuned threshold policy; results must report such cases. Firmware
hooks, telemetry fidelity, and executable-mapping behavior are
platform-specific. Consequently, conclusions are restricted to the tested
runtime commits, devices, memory limits, storage configurations, and
scripted workloads; broader embedded applicability is future validation.

%% =====================================================================
\section{Conclusion}
\label{sec:conclusion}

SAGE-MM specifies how heap configuration, interop allocation, and clean-page
reclamation can be coordinated under a single embedded-process memory budget
through a narrow, safety-guarded online control action, rather than tuned in
isolation. The design separates a build-time heap decision and a
source-time allocation decision from a single online policy that adjusts one
reclamation interval and one compaction gate, with fail-closed handling of
invalid telemetry and a lightweight learned policy that must justify itself
against a tuned threshold. The accompanying protocol makes the resulting
claims falsifiable through a full-factorial ablation, adverse-workload
stress, and endurance runs reported with run-level confidence intervals. We
scope conclusions to the evaluated firmware and workloads and release a
reference kit that reproduces the control logic and safety guards
independently of the proprietary firmware integration.

%% =====================================================================
\section*{Evidence-Availability Statement}
The public repository contains a reference controller, a conservative
mapping helper, an advisory analyzer, and the evaluation protocols. It does
not contain the vendor Mono patch, production access-event integration,
\texttt{Private\_Clean} attribution, raw commercial-device traces, a
complete factorial result bundle, additional-platform results, or the
submitted manuscript source. The kit's file lengths, forced Gen-0 timing,
and fragmentation estimate are proxies and are not used as evidence for
memory, latency, or reliability claims. Simulation-only values in Section~\ref{sec:simulated-results} remain visibly
watermarked and must be removed. Numerical claims remain unset until a
versioned observed-result bundle
supports each with its declared baseline, denominator, statistic, and
confidence interval.

%% =====================================================================
\begin{acks}
The authors thank the anonymous reviewers for feedback that sharpened the
scope, boundary, and evaluation protocol of this work.
\end{acks}

%% =====================================================================
\bibliographystyle{ACM-Reference-Format}
\bibliography{sagemm}

%% =====================================================================
\appendix

\section{Reproducibility and Artifact Manifest}
\label{app:artifact}

To support an ACM artifact evaluation, the released bundle records, where
licensing permits: the vendor Mono runtime name and source revision; the
minimal runtime patch and per-architecture build flags and
\texttt{INITIAL\_ALLOC} values; the exposed host interfaces; the native
reclamation helper; the analyzer rules and accepted diagnostics; the
controller configuration; the workload driver and state-machine transition
seeds; the build, measurement, and bootstrap scripts; the device reset
script; raw per-run telemetry with OOM/kernel logs; and anonymized aggregate
production traces. Each result bundle carries a machine-readable manifest with
the Git revision, device-ID pseudonym, start/end UTC, treatment, seed, and
checksum, so that every number in the paper traces to one versioned source.

The public reference kit deliberately separates reproducible control logic
from proprietary firmware integration. Its user-space compaction event and GC
telemetry are analogues clearly marked as such; they validate control flow,
safety guards, and the statistical pipeline, but are not evidence for the
memory, latency, or reliability claims, which require the production artifact
above.

\section{Notation}
\label{app:notation}

\begin{table}[h]
  \small
  \begin{tabular}{@{}ll@{}}
    \toprule
    \textbf{Symbol} & \textbf{Meaning} \\
    \midrule
    $L_{\mathrm{gc}}$ & GC pause over the preceding interval (ms) \\
    $F_h$ & managed-heap fragmentation ratio \\
    $P_f$ & page-fault rate (minor+major per measured second) \\
    $\Delta M$ & positive resident-memory growth (MiB) \\
    $\mathbf{x}$ & clamped dimensionless feature vector \\
    $y$ & observed normalized pressure (Eq.~\ref{eq:pressure}) \\
    $\hat{y}$ & predicted next-interval pressure (ridge) \\
    $T$ & reclamation interval, bounded to $[T_{\min},T_{\max}]$ \\
    $\eta,\lambda$ & ridge learning rate and regularization \\
    $G,I,R,C$ & heap build, interop, reclamation, controller factors \\
    \bottomrule
  \end{tabular}
\end{table}

\end{document}
